% Generated by tools/edn_latex.py from reports/edn.md.
% Do not edit: change reports/edn.md and run `make edn`.
\ednentry{
  id = {EDN-15},
  title = {UC1 required fields after measuring the fill rates, plus SC-06 for absent optional fields},
  note = {\textbf{Evidence corrected on 2026-09-29.} The first run of \texttt{fillrates.\allowbreak{}py} printed fill rates with one decimal, so \texttt{seller\_\allowbreak{}type}'s 99.986\,\% appeared as 100.0.
The script now prints three decimals and \texttt{fillrates\_\allowbreak{}results.\allowbreak{}txt} was regenerated. The decision is unchanged: 14 missing rows in 97,889 leave it effectively never absent.
The training-scope count moved from 97,913 to 97,889 when the script was brought in line with EDN-22, which drops the listings registered after the reference date.},
  date = {2026-09-29},
  milestone = {M1: Project Inception (requirements and success criteria; shapes M3: Quality Assurance)},
  topic = {Requirements Engineering, API Design, Performance Criteria},
  participants = {@lukas2510},
  decision = {\texttt{/\allowbreak{}predict} requires \texttt{body\_\allowbreak{}type} and \texttt{seller\_\allowbreak{}type} in addition to the core FR-01 asked for before (\texttt{make}, \texttt{model}, \texttt{registration\_\allowbreak{}date}, \texttt{mileage\_\allowbreak{}km\_\allowbreak{}raw}, \texttt{power\_\allowbreak{}kw}, \texttt{fuel\_\allowbreak{}category}, \texttt{transmission}, \texttt{country\_\allowbreak{}code}). The remaining basic-set fields stay optional. FR-01 no longer prescribes how the model handles an absent optional field; it states the observable outcome instead, and the new SC-06 of the problem specification bounds it: with each optional field masked on its own, and with all of them masked at once, the point model's MdAPE stays at or below 1.5 times its full-input MdAPE. This amends EDN-06, which set SC-01 to SC-05; because NFR-01 gates deployment on all of them, the gate now also covers partial UC1 inputs. Whether the rarely-missing optional fields need the random masking planned for the UC2 interval models is left to M2/M3, when the pipeline exists and the effect can be measured.},
  rationale = {\emph{Alternatives considered:}
\begin{itemize}[leftmargin=*, topsep=2pt]
\item \textbf{Option A: keep all eight basic-set fields optional and rely on LightGBM's native missing-value handling.}
\newline Pros: the smallest required set, so UC1 stays easy to call and clearly different from UC2; no training-time work at all.
\newline Cons: measured to be unsupported for two fields. \texttt{body\_\allowbreak{}type} is filled in 100.000\,\% of the training listings and \texttt{seller\_\allowbreak{}type} in 99.986\,\% (14 of 97,889 rows), so the model effectively never sees them absent and a request that omits them is sent to the default side of every split on that feature with no training example backing that direction. The contract would promise behaviour the training data cannot show.
\item \textbf{Option B: keep all eight optional and train the point model with the same random masking as the UC2 interval models.}
\newline Pros: friendliest contract; one masking scheme for both model families, which also keeps the point estimate and the interval bounds from drifting apart on partial inputs (FR-07 widens the bounds when they do).
\newline Cons: makes an API requirement depend on a pipeline that does not exist yet (no \texttt{dvc.\allowbreak{}yaml}, issue \#10 still open), so PR \#19 would stay blocked on a modelling decision that cannot be measured before the first report. Masking also costs some accuracy on full inputs, which are the common case for UC1.
\item \textbf{Option C: require all basic features.}
\newline Pros: best accuracy, no missing-value path at serving time at all.
\newline Cons: hostile to users. \texttt{nr\_\allowbreak{}prev\_\allowbreak{}owners} is filled in 61.0\,\% of the listings, \texttt{gears} in 62.5\,\% and \texttt{drive\_\allowbreak{}train} in 76.4\,\%, so even the sellers behind the listings often do not state them.
\item \textbf{Option D (chosen): require the two always-filled fields, keep the rest optional, and bound the outcome with SC-06.}
\newline Pros: closes the measured gap at zero cost to the user, because whoever owns the car knows its body type and, for UC1, is the seller. Separates the API contract (decidable now) from the training mechanism (decidable in M2/M3), so the requirement can be merged without guessing. SC-06 turns FR-01's previously unquantified ``per-field masked-input accuracy check'' into a threshold and puts it in the deployment gate.
\newline Cons: the required set grows from eight fields to ten. SC-06 amends EDN-06, which was already agreed, and it is relative rather than absolute, so it cannot rule out a model that is uniformly mediocre; SC-01 to SC-04 do that.
\item \textbf{Option E: adopt D, but report the masked-input error only, without gating on it.}
\newline Pros: no amendment to EDN-06; the number is still visible in MLflow and in the report.
\newline Cons: leaves NFR-01's gate blind to exactly the case FR-01 allows the client to trigger, which is how the gap arose in the first place.
\end{itemize}
\par
\emph{Why:} FR-01 carried an \texttt{[open]} marker because it mixed two questions: which fields the client must send, which is an API contract, and how the model is trained to honour the optional ones, which is a modelling mechanism. The fill rates that the second question depended on had never been measured. Measured on the scoped, deduplicated training data (97,889 listings, \texttt{ES} held out, 11 supported makes), the risk turns out to sit in exactly two fields: \texttt{body\_\allowbreak{}type} and \texttt{seller\_\allowbreak{}type} at 100\,\%, followed by \texttt{nr\_\allowbreak{}doors} (98.8\,\%), \texttt{nr\_\allowbreak{}seats} (97.0\,\%) and \texttt{cylinders\_\allowbreak{}volume\_\allowbreak{}cc} (91.1\,\%), while \texttt{nr\_\allowbreak{}prev\_\allowbreak{}owners} (61.0\,\%), \texttt{gears} (62.5\,\%) and \texttt{drive\_\allowbreak{}train} (76.4\,\%) have plenty of natural missingness. Only 30.7\,\% of training rows have every one of the six remaining optional fields present and 469 rows (0.5\,\%) have none of them, so the model sees partial rows including the extreme case, just never a row without \texttt{body\_\allowbreak{}type} or \texttt{seller\_\allowbreak{}type}. Requiring the two is therefore the cheapest correct fix, and it costs the user nothing. The rest is left to the modelling plan on purpose: the mechanism cannot be chosen on evidence before the pipeline runs, and blocking the requirements page on it would have held up the M1 deliverable for a decision that belongs to M2/M3. SC-06 is what makes that safe, and it is relative to the model's own full-input MdAPE, like SC-03 is relative to the baseline, because no reference value exists yet and an absolute threshold would have been invented rather than derived.},
  ai = {Information seeking, Alternative generation, Alternative assessment, Recommendation},
  response = {Accepted with modifications},
  assessment = {Asked which decisions had to be settled before PR \#19 could be merged, AI named FR-01's open item first and called it a merge blocker. Asked to re-examine that, it downloaded the raw dataset from the Zenodo DOI, reproduced the scope, deduplication and make filter (which matched the brief's 11 makes as a cross-check), measured the per-field fill rates, and then corrected two of its own claims: the item does not block code that exists, because there is no pipeline yet, and FR-01 already promised a masked-input check, so the real gap was that the check had no threshold and no link to NFR-01's gate. It also reframed the requirement as mixing contract and mechanism, which is what made a clean split possible. Lukas accepted the recommendation and the field list; the modification is that the mechanism question stays open on purpose instead of being decided together with the contract.},
  aievidence = {Claude Code session on 2026-09-29: Lukas asked (translated from German) ``which things urgently still need deciding in requirements.md before we can merge, and why?'', then ``analyse that again, then explain it to me and give me a reasoned recommendation on what the best way is''; AI measured the fill rates, revised its own analysis, proposed options A to E and recommended D with SC-06 in the gate; Lukas replied ``I like your recommendation, please change it that way in the PR''.},
  evidence = {\href{https://github.com/mlops-2627q1-mds-upc/MLOPS_Recommenditos/blob/main/docs/docs/specification.md}{Specification} FR-01; \href{https://github.com/mlops-2627q1-mds-upc/MLOPS_Recommenditos/blob/main/docs/docs/problem-spec.md}{problem specification} SC-06; \href{https://github.com/mlops-2627q1-mds-upc/MLOPS_Recommenditos/blob/main/reports/analysis}{reports/analysis/} (\texttt{fillrates.\allowbreak{}py}, run of 2026-09-29); \href{https://github.com/mlops-2627q1-mds-upc/MLOPS_Recommenditos/blob/main/docs/docs/project-brief.md}{project brief} section 4; \hyperref[EDN-06]{EDN-06}; \href{https://github.com/mlops-2627q1-mds-upc/MLOPS_Recommenditos/pull/19}{PR \#19}.},
}
